%%%%%%%%%%%%%%%%%%%%%%%%%%%%%%%%%%%%%%%%%%%%%%%%%%%%%%%%%%%%%%
% REVISED DISCUSSION — drop-in replacement for lines 276–310 of main.tex
% Streamlined per Shinichi's comments:
%   (1) merged §4.1 with the conditional/heterogeneity material from §4.2;
%   (2) cut result-level numbers (kept in Results);
%   (3) combined the bias, decline, and reporting-standards paragraphs.
% Opens by answering the two research questions (Q1 survives / Q2 attenuated).
% ~1,200 words / 7 paragraphs (was ~1,920 words / 12 paragraphs).
% All original \cite/\textcite keys preserved.
%%%%%%%%%%%%%%%%%%%%%%%%%%%%%%%%%%%%%%%%%%%%%%%%%%%%%%%%%%%%%%
\section{Discussion}

By updating the datasets from the two original meta-analyses and reanalysing them within a unified, phylogenetically informed multilevel framework, we can now answer the two questions we set out with. First, the qualitative conclusion of both foundational syntheses survives: across an expanded, multilingual, and grey-literature--inclusive search, observers experimentally exposed to social information were on average more likely to favour the associated mate, and the pooled effect remained positive and statistically significant on both the odds-ratio and standardised-mean-difference scales. Second, the effect is smaller than previously implied---roughly half the original magnitude when the new evidence is considered on its own, with confidence intervals that span zero, and, under the more conservative bias correction, an adjusted mean close to zero. Mate-choice copying is therefore real, but more modest and more context-dependent than \textcite{jonesMechanismsSocialInfluence2019} and \textcite{daviesMetaanalysisFactorsInfluencing2020} implied---a pattern that is itself informative and that we explore below.

%%%%%%%%%%%%%%%%%%%%%%%%%%%%%%%%%%%%%%%%%%%%%%%%%%%%%%%%%%%%%%
\subsection{Mate-choice copying as social learning, and its consequences for sexual selection}

Theory frames mate-choice copying as one expression of a general capacity to use public information, alongside socially guided foraging, habitat selection, and predator avoidance~\cite{danchinCulturalFliesConformist2018, galefjr.SocialLearningAnimals2005, valonePublicInformationAssessment2002}, with a central prediction that such information \emph{supplements}, rather than \emph{overrides}, an observer's own assessment of a prospective mate~\cite{vakirtzisMateChoiceCopying2011, ueharaMatechoiceCopyingBayesian2005}. A reliably positive yet modest average effect is exactly what that prediction implies, and it reframes the biological question from \textit{whether} animals copy to \textit{when} they do so. Theory does not predict indiscriminate copying: social-learning strategies such as ``copy when uncertain'' and ``copy when young or inexperienced'' specify the narrow conditions under which relying on others outperforms personal information~\cite{lalandSocialLearningStrategies2004, galefjr.SocialLearningAnimals2005}. Seen this way, the recent, well-powered studies that report weak, null, or context-restricted copying~\cite{nobelNoEvidenceMate2023, belkinaMateChoiceCopying2021, pusiakRelativeSexualAttractiveness2020} are not evidence against the phenomenon but evidence that, as theory predicts, it is expressed selectively---animals discount or reverse the behaviour when public information is unreliable, when demonstrators are sperm-limited, or when individual differences in personality intervene~\cite{loyauWhenNotCopy2012, dugatkinDisruptionHypothesisDoes2003, nobelImportancePopulationHeterogeneities2022}.

This conditional view also reframes the most prominent feature of our results---large and largely unexplained heterogeneity---as biological signal rather than statistical noise. Variation accumulated within studies and among species, rather than between laboratories or across the phylogeny, which is what we would expect if the expression of copying is tuned to the observer's experience, the reliability and consistency of the social information, and the local ecological context. \textcite{jonesMechanismsSocialInfluence2019} found exactly such conditionality, with stronger copying among inexperienced (virgin) females; our updated data, however, did not reproduce this virginity effect and, in the new studies, tended to reverse it. The candidate moderators we could examine---copying mechanism, cue modality, observer mating status, and demonstrator age---were explored only in post-hoc, non-pre-registered analyses and explained no detectable share of the heterogeneity, consistent with the relevant variation residing at finer or more interactive scales than our coding captured. That copying occurs even in animals lacking complex sociality, most notably \textit{Drosophila}, further suggests that it taps domain-general learning mechanisms shared widely across animals rather than lineage-specific cognition~\cite{bridgesAnimalBehaviourConformity2019a, danchinCulturalFliesConformist2018}---an interpretation consistent with the negligible phylogenetic signal we recovered.

One distinction among our exploratory moderators carries particular weight for these evolutionary consequences: whether observers copy in a \textit{generalized} or an \textit{individual} fashion. Only generalized copying---the acquisition of a preference for a male phenotype or trait rather than for one specific individual---can propagate through a population, bias the targets of sexual selection, and seed the cultural transmission of mating preferences, although theory disagrees on how strongly copying ultimately accelerates or constrains sexual selection~\cite{varelaRoleMatechoiceCopying2018}; individual copying, by contrast, ends with the demonstrated male. The evolutionary reading of mate-choice copying therefore rests on the generalized, transmissible form being robustly expressed, and our data support that premise: generalized copying was statistically indistinguishable from individual copying, and if anything marginally stronger in the new extraction. This makes a mechanism-explicit, pre-registered test of the generalized--individual distinction a high priority for future work.

%%%%%%%%%%%%%%%%%%%%%%%%%%%%%%%%%%%%%%%%%%%%%%%%%%%%%%%%%%%%%%
\subsection{Knowledge gaps and future opportunities}

Three features of the evidence base temper these conclusions and set the agenda for future work. First, we found strong small-study effects on both scales; under the bias-corrected intercept~\cite{nakagawaMethodsTestingPublication2022} the adjusted mean collapsed towards---and for Hedges' \textit{g} marginally below---zero, whereas the bias-robust estimator~\cite{yangRobustPointVariance2024} retained a positive but reduced signal. The honest interpretation is that the true effect very likely lies below the naïve pooled estimate, and that a much smaller effect cannot be excluded on present evidence. Second, for the binary outcomes the effect-size magnitude declined with publication year, echoing the familiar pattern in which early, large effects are gradually tempered by more rigorous and better-powered designs~\cite{jennionsRelationshipsFadeTime2002, korichevaTemporalInstabilityEvidence2019}; several of the most recent studies in our sample report weak or null copying~\cite{belkinaMateChoiceCopying2021, cirinoRobustMatePreferences2023, pusiakRelativeSexualAttractiveness2020, nobelNoEvidenceMate2023}. Third, primary studies quantify copying in incompatible ways, and the two source meta-analyses reported their results in different metrics, forcing approximate conversions and leaving a handful of studies whose values disagreed materially between datasets. Although our conclusions proved robust to how these contentious values were handled, the broader lesson is methodological: the field would benefit from registered reports, the routine publication of null and small-sample results, and minimum reporting standards---group means, standard deviations, sample sizes, and raw choice counts---so that future updates rest on less-censored, directly comparable data~\cite{daviesMetaanalysisFactorsInfluencing2020}.

A further gap concerns taxonomic breadth. Although our combined sample spanned 33 species across arthropods, fish, and other vertebrates, the evidence remains concentrated in two model genera: \textit{Drosophila} and \textit{Poecilia} together contributed roughly two-thirds of all effect sizes. Reassuringly, the uni-moderator analyses returned positive means in all three taxonomic groups on both scales, with only weak support for among-group differences, so we treat the apparent gradient cautiously. Amphibians, reptiles, non-poeciliid fishes, and a broader range of birds and mammals remain effectively absent. Expanding the phylogenetic spread of the evidence would both test the generality of mate-choice copying and sharpen phylogenetic comparative tests that, at present, have little signal to work with.

%%%%%%%%%%%%%%%%%%%%%%%%%%%%%%%%%%%%%%%%%%%%%%%%%%%%%%%%%%%%%%
\subsection{Conclusion}

Mate-choice copying is a real and taxonomically widespread behaviour: after incorporating seven additional years of evidence and analysing it with contemporary phylogenetic multilevel methods, we still recover a positive and significant overall effect on both the odds-ratio and standardised-mean-difference scales. At the same time, the updated picture is more sober than the original meta-analyses suggested. The pooled effect is moderate at best, is surrounded by substantial heterogeneity that our models could not attribute to study, taxon, or phylogeny, and bears clear fingerprints of small-study and selective-reporting bias that, under the more conservative correction, draw the estimate close to zero. The qualitative claims of \textcite{jonesMechanismsSocialInfluence2019} and \textcite{daviesMetaanalysisFactorsInfluencing2020} hold, but their estimated magnitudes were probably optimistic. Realising the promise of mate-choice copying as a driver of sexual selection and cultural evolution will require broader taxonomic sampling, mechanism-focused and pre-registered tests of the moderators that potentially influence its considerable heterogeneity, and a publication culture that faithfully reports null findings. Finally, we offer this update as a template for how meta-analyses in ecology and evolution can be revisited---combining new evidence with the original data, harmonising incompatible effect-size metrics, and stress-testing conclusions against publication bias and contentious extractions.
